% Discussion --- defensible claims and honest limitations. This section
% is the guard against over-claiming; it states explicitly what we do NOT
% claim.
\section{Discusión}
\label{sec:discussion}

\subsection{Afirmaciones defendibles}
Nuestros resultados sustentan cuatro afirmaciones, planteadas de forma
conservadora.
\begin{enumerate}[leftmargin=*]
    \item Los transformadores temporales con una rama fenológica son los
    perceptores individuales más fuertes sobre PASTIS-R (TSViT-pheno,
    mIoU 0.6253), y el stacking heterogéneo sobre miembros
    decorrelacionados alcanza un F1 macro a nivel de parcela en el fold de
    prueba de 0.7470 (exactitud 0.8490), mejorando el stacking de tres
    miembros (de 0.6359 a 0.6477 out-of-fold) al arbitrar clases raras.
    El perceptor de producción desplegado es el más simple Voting-3 (F1
    macro a la par, mejor en regiones nuevas); usado como perceptor del
    agente eleva la exactitud a nivel de parcela de 0.8539 a 0.9375
    ($+8.4$ puntos) sobre 14{,}688 parcelas reales. Esto confirma que la
    separación Be My Eyes rinde frutos precisamente porque el razonador
    congelado hereda la exactitud del perceptor desplegado.
    \item Una rama contrastiva visión-lenguaje~\cite{li2025farslip}
    aporta una señal complementaria útil pero modesta (la ganancia del
    stacking de cinco miembros sobre tres es $+0.0118$ de F1 macro). Es
    un complemento, no un clasificador autónomo.
    \item El patrón \emph{Be My Eyes}~\cite{huang2025bemyeyes} produce
    una capa conversacional anclada en la que cada cifra citada es
    trazable a un resultado de herramienta auditable, con un razonador
    congelado intercambiable entre la nube (Gemini~3.5~Flash) y on-premises
    (Qwen3-30B-A3B) para la soberanía de datos.
    \item La transferencia espacial es limitada por región, no zero-shot:
    de Francia a Cataluña recupera mIoU 0.2468 solo tras ajuste fino
    few-shot, e incluso una transferencia dentro del mismo país de Francia
    a Bretaña cae a F1 0.21, en línea con los límites de transferibilidad
    reportados~\cite{harvesting2026alphaearth}. Aun así, la transferencia
    \emph{sí} funciona donde la fenología es separable --- el campeón
    entrenado en Francia alcanza F1 macro 0.266 en Baja Sajonia (DE4),
    rescatando los cereales de invierno --- y el modelo de fundación
    reduce el presupuesto de etiquetas few-shot en aproximadamente un
    orden de magnitud (AlphaEarth supera a Sentinel-2 crudo en $+0.111$ de
    F1 con $k{=}1$ en EuroCropsML LV$\rightarrow$EE).
    \item Agrupar regiones en un único clasificador multirregión amplía la
    taxonomía de 18 a 30 hojas sin degradar la lectura gruesa, pero no
    rescata \emph{ninguna} clase fina por encima del umbral de producción
    de 0.85: el cuello de botella es la separabilidad fenológica en un
    embedding anual, no el número de regiones. Reportamos este resultado
    negativo honesto para acotar la afirmación.
\end{enumerate}

\subsection{Lo que explícitamente no afirmamos}
Para evitar sobreafirmar, delimitamos las fronteras de nuestra
evidencia.
\begin{itemize}[leftmargin=*]
    \item \emph{No} afirmamos que un VLM ajustado supere al razonador
    congelado Gemini. La capa de lenguaje es un componente de
    comunicación, explicación y orquestación de herramientas, no un
    clasificador de píxeles; su valor es el anclaje, la auditabilidad y
    la opción de soberanía on-premises.
    \item \emph{No} afirmamos un umbral de clasificación
    (p.\ ej.\ F1~$\geq$~0.80) en los trópicos. La transferencia
    Europa$\rightarrow$Brasil/India con WorldCereal se mide con etiquetas
    reales de verdad de terreno y no supera F1~0.72 ni siquiera tras
    adaptación few-shot; el maíz zero-shot colapsa a F1~0.0095.
    \item \emph{No} afirmamos que la transferencia zero-shot funcione
    fuera de Francia. El experimento Francia$\rightarrow$Cataluña muestra
    el fallo zero-shot y la recuperación few-shot; una curva de
    $k$-shots sobre EuroCropsML~\cite{reuss2025eurocropsml} caracteriza
    el presupuesto de datos requerido.
    \item \emph{No} afirmamos una mIoU de TSViT superior a 0.75. PASTIS-R
    se satura cerca de 0.70 mIoU para este conjunto de clases, y la rama
    fenológica aporta del orden de $0.3$ puntos, no varios.
\end{itemize}

\subsection{Limitaciones}
Permanecen tres limitaciones estructurales. Primero, seis clases
minoritarias con fenología ambigua arrastran a la baja el F1 macro
global; abordarlas requiere datos adicionales o aumentación dirigida, no
más modelos del mismo tipo. Segundo, la cobertura nubosa puede ocultar
una gran fracción de píxeles en las estaciones lluviosas, lo que motiva
la fusión con el radar Sentinel-1. Tercero, el benchmark de la capa
conversacional sobre AgroMind~\cite{ruan2025agromind} está pendiente de
la evaluación en H100 (US-069); reportamos el protocolo y la estructura
de columnas pero ninguna cifra sintética.

\subsection{Eficiencia de costos (FinOps)}
El sistema está diseñado para operar a bajo costo operativo. Los cuatro
componentes de servicio (API, frontend, teselado, worker de inferencia)
corren en Cloud Run con escalado a cero, y el razonador on-premises H100
es una instancia spot de Azure activada solo bajo demanda
(aproximadamente tres horas al día). El costo operativo en estado
estacionario es de aproximadamente US\$115 al mes, mientras que el
presupuesto único de entrenamiento --- las ventanas spot H100 y L4 usadas
para ajustar los segmentadores y la destilación de FarSLIP --- totaliza
US\$262--602. El diseño de razonador congelado es en sí mismo una palanca
de costo: como el LLM nunca se ajusta finamente como clasificador, no hay
un ajuste fino costoso en GPU de un modelo frontera en la ruta crítica, y
el cómputo pesado se confina al entrenamiento único de la pila de
percepción.

\subsection{Reproducibilidad}
Cada figura de la Sección~\ref{sec:results} está anclada a un artefacto
versionado mediante comentarios \texttt{\% src:} en el código fuente
LaTeX; los conjuntos de datos y los pesos se rastrean con DVC y las
corridas de experimentos con MLflow bajo etiquetas
\texttt{data\_version} y \texttt{code\_version}. El repositorio se
publica bajo Apache~2.0.
